% ═══════════════════════════════════════════════════════════════════════
%  Stratégie 2 --- XAUUSD niveaux x10. Tous les chiffres viennent des
%  macros de tables/x10_headline.tex.
% ═══════════════════════════════════════════════════════════════════════
\section{Méthodologie}
\label{sec:x10_methodologie}

\lettrine[lines=2, findent=2pt, nindent=0pt]{L}{a} question que pose ce
chapitre n'est pas « comment la stratégie a-t-elle été mesurée~? » mais
« qu'est-ce qui empêchait le mesureur de se tromper en sa faveur~? ». Un
backtest se laisse toujours convaincre quand on le relance assez souvent. Le
dispositif décrit ici est celui qui rend ce rapport publiable, y compris ---
et surtout --- avec un verdict négatif.

\subsection{La spécification avant le code}

La totalité de la mécanique a été écrite avant la première ligne de code~:
horloge et bornes de séance, définition des niveaux, formule de l'\code{ATR}
et son amorçage, cinématique, construction du contexte, machine à états,
règles d'exécution, dimensionnement, modèle de coût, contrat de trace. Douze
points laissés ambigus par le brief ont été tranchés à ce moment-là, pas
pendant la mesure.

L'intérêt n'est pas documentaire. Quand la spécification précède le code, un
comportement gênant découvert plus tard est un \emph{défaut du moteur}, et il
se corrige. Quand elle le suit, le même comportement devient une
« particularité de l'implémentation », et il se garde s'il arrange les
chiffres.

\subsection{Données, séances et fuseaux}

Le graphique de décision est le \code{M5}~; la résolution des ordres se fait
sur les barres d'une minute qui le composent. L'échantillon de sélection court
du 1\textsuperscript{er} janvier 2019 au 31 décembre 2025, soit
\XTenSampleBarsMOne{} barres d'une minute, \XTenSampleBarsMFive{} barres de cinq
minutes et \XTenSampleSessions{} séances.

La séance de l'or court du dimanche 18:00 au vendredi 17:00, heure de New
York. C'est cette horloge --- et non l'heure UTC ni le fuseau du courtier ---
qui découpe les séances, réinitialise le \code{VWAP} et déclenche la clôture
de fin de journée. Les trois moteurs doivent en reproduire la convention, le
moteur MetaTrader~5 par une conversion explicite de l'heure serveur, vérifiée
sur les quatre bascules d'heure d'été.

\subsection{Le décalage des séries H1 --- la règle anti-anticipation}

\begin{warningbox}
\textbf{C'est le point où une stratégie multi-échelles se trompe, et le dépôt
l'a déjà payé une fois.} Agréger des barres d'une minute en barres horaires
produit une barre étiquetée \code{08:00} qui contient la clôture de
\code{08:59}. Reportée telle quelle sur le graphique de cinq minutes, elle
rend visible à 08:00 une information qui n'existera qu'une heure plus tard. Un
travail antérieur du dépôt affichait un Sharpe de recherche flatteur sur trois
configurations~; après correction, les trois sont devenues négatives. L'alpha
était entièrement un artefact.
\end{warningbox}

La spécification impose donc une chaîne non négociable pour toute série
horaire --- moyenne mobile exponentielle à 50~périodes, \code{ATR} horaire,
panier dollar~: agréger, \emph{décaler d'une barre}, puis seulement reporter
sur la grille de cinq minutes. Un test dédié doit rougir sur une version sans
décalage.

Le \code{VWAP} de séance, lui, n'est pas une série horaire et n'obéit pas à
cette règle~: il se calcule sur les barres de cinq minutes elles-mêmes et est
entièrement connu à la clôture de la barre de décision. Confondre les deux
familles est la première erreur possible sur cette stratégie~; la
section~\ref{sec:x10_mecanique} les sépare explicitement.

\subsection{Correction de la convention de datation}

La première campagne traitait l'horodatage du parquet comme le début de la
minute. La réconciliation a montré qu'il s'agissait de sa clôture, selon la
convention d'export de LEAN~: la grille de cinq minutes de la référence était
donc décalée d'une minute par rapport à la spécification et aux moteurs
d'exécution. La correction redéfinit d'abord chaque ligne comme le début de sa
minute, puis agrège en M5.

La campagne a été rejouée sans changer un seuil, une règle, une configuration
ou une clé du registre d'essais. Ce rejeu corrige la mesure~; il ne constitue
ni une nouvelle optimisation ni une re-sélection. La version antérieure est
archivée et la comparaison complète des neuf critères est publiée ci-dessous.

% GÉNÉRÉ par scripts/build_x10_report_assets.py — ne pas éditer à la main.
% Toute valeur vient de results/xau_x10/*.json.
\begin{table}[htbp]
\centering
\footnotesize
\caption{Les neuf critères avant et après la correction de la convention de datation des minutes. Un seul change de camp, et il change en faveur de la stratégie~; le verdict, lui, ne bouge pas. Ces deux archives portent sur la campagne Python : la mesure MT5 ajoutée ensuite figure dans la table de décision consolidée.}
\label{tab:x10_v1_v2}
\begin{tabular}{c p{4.6cm} r c r c}
\toprule
\textbf{\#} & \textbf{Critère} & \multicolumn{2}{c}{\textbf{Avant}} & \multicolumn{2}{c}{\textbf{Après}} \\
\midrule
1 & Trades in-sample 2019 $\rightarrow$ 2025-12-31 & $1\,723$ & \textcolor{rsiGreen}{\textbf{GO}} & $1\,762$ & \textcolor{rsiGreen}{\textbf{GO}} \\
2 & Espérance nette / profit factor & $-0{,}1602$ R & \textcolor{combinedBurgundy}{\textbf{échec}} & $-0{,}1679$ R & \textcolor{combinedBurgundy}{\textbf{échec}} \\
3 & DSR (déflaté par distinct\_trials("xau\_x10")) & $0{,}000$ & \textcolor{combinedBurgundy}{\textbf{échec}} & $0{,}000$ & \textcolor{combinedBurgundy}{\textbf{échec}} \\
4 & PBO (CSCV, matrice 27 $\times$ jours) & $0{,}4155$ & \textcolor{combinedBurgundy}{\textbf{échec}} & $0{,}2460$ & \textcolor{rsiGreen}{\textbf{GO}} \\
5 & Plateau & $0$ sur 27 & \textcolor{combinedBurgundy}{\textbf{échec}} & $0$ sur 27 & \textcolor{combinedBurgundy}{\textbf{échec}} \\
6 & Walk-forward annuel & $0$\,\% & \textcolor{combinedBurgundy}{\textbf{échec}} & $0$\,\% & \textcolor{combinedBurgundy}{\textbf{échec}} \\
7 & Spread $\times 1{,}5$ & $-0{,}2012$ R & \textcolor{combinedBurgundy}{\textbf{échec}} & $-0{,}2078$ R & \textcolor{combinedBurgundy}{\textbf{échec}} \\
8 & MT5 modèle 1, 2022-11-04 $\rightarrow$ 2025-12-31 & \emph{non mesuré} & \textcolor{combinedBurgundy}{\textbf{échec}} & \emph{non mesuré} & \textcolor{combinedBurgundy}{\textbf{échec}} \\
9 & OOS 2026, lecture unique & \emph{non mesuré} & \textcolor{combinedBurgundy}{\textbf{échec}} & \emph{non mesuré} & \textcolor{combinedBurgundy}{\textbf{échec}} \\
\bottomrule
\end{tabular}
\end{table}


Un seul critère change de camp, le \code{PBO}, dans un sens favorable à la
stratégie. La décision de ne pas déployer demeure inchangée.

\subsection{Découpage in-sample / hors échantillon}

Tout ce qui est postérieur au 31 décembre 2025 est \textbf{verrouillé}. Le
verrou n'est pas une intention~: c'est un contrôle exécuté à chaque étape de
la campagne, sur les trois index effectivement consommés --- minute, cinq
minutes, séance. La campagne publie la date maximale de chacun comme preuve.

Une conséquence mérite d'être dite. Les barres du soir du 31 décembre 2025
appartiennent à une séance étiquetée 2026~: elles ont été retirées, plutôt que
de laisser une étiquette de la tranche gelée entrer dans l'index de sélection.
Le choix est conservateur et sans effet sur les conclusions.

\subsection{Budget d'essais et registre de sélection}

Chaque configuration évaluée consomme un \emph{essai}, et les essais se
comptent. Le budget de la famille a été fixé à \XTenTrialsCeiling{} avant la
première mesure~: vingt-sept pour la grille, sept pour les ablations, le reste
en réserve. La campagne en a consommé \XTenTrialsCount, laissant
\XTenTrialsReserve{} en réserve intacte.

\begin{insightbox}
\textbf{Le budget d'essais n'est pas une discipline morale, c'est un
paramètre de calcul.} Le \emph{Deflated Sharpe Ratio} publié en
section~\ref{sec:x10_robustesse} déflate le Sharpe observé par le nombre
d'essais distincts. Avec \XTenTrialsCount{} essais, le Sharpe attendu du
\emph{meilleur} d'entre eux vaut \metric{\XTenExpectedMaxSharpe} par pur
hasard --- c'est le seuil que la stratégie devait battre, et non zéro.
Dépasser le budget n'est pas une faute de procédure~: c'est une invalidation
arithmétique du chiffre publié.
\end{insightbox}

\subsection{La table de décision, gelée avant la mesure}

Neuf critères, neuf seuils, trois verdicts possibles, écrits et committés
avant que le moindre résultat soit lu. Deux clauses de cette table comptent
plus que les autres.

\begin{description}
  \item[Un critère non mesuré vaut échec.] Jamais « non applicable ». Cette
  clause a un coût immédiat dans ce rapport~: le critère OOS 2026 est en
  échec par non-mesure, et il est compté comme tel
  (section~\ref{sec:x10_robustesse}).
  \item[Le plateau, jamais le pic.] La configuration retenue doit
  être entourée de voisins qui tiennent, pas être le meilleur point d'une
  carte. Et si aucune configuration ne passe ce test, la table prévoit
  explicitement qu'\emph{il n'y a pas de configuration retenue} et que le
  verdict est négatif. C'est exactement ce qui s'est produit.
\end{description}

\subsection{Le modèle de coût}

Le spread de référence retenu pour la sélection est constant, à
\XTenSpreadUsd{} l'aller-retour, relevé au catalogue du courtier. Ce choix a
été gelé avant la campagne, avec une sensibilité obligatoire au doublement
du coût, en plus de celle à une fois et demie. Le dump du courtier obtenu
ensuite permet une sensibilité supplémentaire au spread mesuré~; elle est
publiée sans rouvrir la grille et sans remplacer le modèle de sélection.

Un spread constant est trop cher en séance liquide et trop bon marché en
heures creuses~; il pèse aussi mécaniquement de moins en moins à mesure que le
prix de l'or monte. C'est une limite à publier, pas un motif de rejouer la
grille~: la section~\ref{sec:x10_couts} en tire toutes les conséquences, y
compris celle qui dérange.

\subsection{Les trois moteurs : vectoriel, QuantConnect, MetaTrader~5}

La même spécification est portée sur trois moteurs qui ne servent pas à la
même chose. Le moteur vectoriel explore --- il évalue une grille entière en
quelques secondes. Le moteur QuantConnect rejoue la stratégie
\emph{événement par événement}, comme le ferait une exécution réelle. Le
moteur MetaTrader~5 est celui qui passerait les ordres, sur le flux du
courtier, qui n'est pas celui des deux autres.

Les trois ne peuvent pas donner le même chiffre, et viser l'égalité serait le
signe qu'on a idéalisé l'exécution. La section~\ref{sec:x10_reconciliation}
décrit l'échelle de lecture à cinq barreaux qui sert à localiser un écart
plutôt qu'à le nier.

\subsection{La revue à contexte vierge}

\begin{examplebox}[title=\textbf{\color{white}\textsf{Ce qu'une relecture indépendante a trouvé}}]
Avant toute mesure, le moteur de référence a été relu par un examinateur
disposant de la spécification et du code, mais d'aucun élément de la
discussion qui les avait produits. Cette relecture a trouvé un défaut réel~:
la clôture de séance était testée comme un \emph{seuil} (« il est plus de
16:55~») et non comme une \emph{fenêtre} (« l'heure tombe entre 16:55 et
18:00~»). Conséquence~: toute position ouverte le soir, à partir de 18:15,
aurait été liquidée à la barre suivante, puisque son heure locale est
supérieure à 16:55.

Le défaut a été corrigé \textbf{avant} que la moindre performance soit lue. Ce
n'est pas une anecdote~: c'est la garantie que les chiffres de ce rapport
portent sur la stratégie spécifiée, et non sur une version mutilée de
celle-ci. Un défaut de cette nature découvert \emph{après} la lecture aurait
rendu la campagne non concluante, pas corrigible.
\end{examplebox}

\subsection{Reproductibilité}

Toutes les mesures publiées ici sortent d'une seule commande, à graine fixe,
dont les sorties sont des fichiers de résultats versionnés. Le rapport ne
recopie aucun chiffre à la main~: chaque nombre de la prose et chaque cellule
de tableau est produite par un générateur unique qui lit ces fichiers, et un
test de la suite du dépôt échoue si un chiffre publié cesse de correspondre à
sa source. La révision du dépôt qui a produit les résultats est
\code{\XTenGitHead}.
